\scanpage{110}
% L1--4
\uline{Corollaire} 1 Supposons que $X$, \deleted{\unsure{soient des}} \inserted{$Y$ soient} des espaces (T$_0$), \deleted{schématiques \unsure{où|qui}} admettant des bases d'ouverts formées d'ens \deleted{quasi-compacts} rétrocompacts et qu. cpts pour la top. constructible [donc compacts pour la dite], \deleted{et que} \inserted{Conditions équivalentes :}

% L5
a) $f:X\longrightarrow Y$ \deleted{soit une telle que} \inserted{est} quasi compacte

% L6--9
\deleted{[par exemple $X$ Alors} b) $f^{\mathrm{cons}}$ est une \deleted{\unsure{application}} application propre d'espaces loc. cpts, \deleted{\unsure{ceci signifie|i.e.} $f$ [\unsure{et ses}] \unsure{parties de} $f$} \deleted{[et \inserted{alors $f$} transforme parties \unsure{antitypiques|antypiques|atypiques|analytiques} en parties antitypiques}

\inserted{c) $f$ transforme parties antitypiques en parties antitypiques, et les fibres de $f$ sont quasi-compactes [\unsure{appliq. de|appl. :} Bourbaki, Top \deleted{\unsure{II}} I, 10, 2, th. \uline{1}]}
% slanted

% L9--12
[N.B. $f$ en général ne transforme pas parties constructibles en parties constructibles, \deleted{et bien mieux on peut faire varier par \unsure{ouverts, idéal} ne transforme pas} \inserted{cf. plus bas]}

% L13--15
\uline{Corollaire 2} Si les conditions précédentes \inserted{sont vérifiées,} supposons $=$ $Y$ donc $X$ qu. cpt, donc $X^{\mathrm{cons}}$ \unsure{et} $Y^{\mathrm{cons}}$ \uline{compacts}. Conditions équivalentes

% L16--17
(a) $f$ transforme parties constructibles en parties constructibles [i.e. transforme ``\unsure{ouvert et fermé|ouverts et fermés}'' en ``ouvert et fermé'']

% L18--19
(b) $f$ transforme parties constructibles en parties réunions de constructibles [i.e. ``ouvert et fermé'' en ``ouvert'']

% L20--21
(c) $f^{\mathrm{cons}}$ est une application \uline{ouverte}. \deleted{\unsure{et}} \deleted{\unsure{(resp. ce qui est)}}

Ces conditions sont vérifiées si $f$ est \uline{bijective}.

% L22--24
\uline{Cor.} 3 Si les conditions du lemme 1 sont vérifiées, alors (i) $f(X)$ est antitypique (ii) plus gén, Une partie \inserted{$E$} de $f(X)$ est antitypique \uline{ssi} $f^{-1}(E)$ l'est (iii) Si $f$ est surjective, une partie de $Y$ est constructible ssi $f^{-1}(Y)$ l'est.
% archive sheet number: 109
